\section{Viscoelastic three-dimensional bubble}
\label{sec:viscoelastic}

\status{\textbf{Closed on the uniform mesh}: all twelve runs of the matrix
(six $(\lambda_1,\beta)$ combinations $\times$ two methods) completed 2001 steps
at $25\times25\times75$ with zero non-finite values. The reference resolution is
complete for the cusp case in VOF and running in front-tracking. The adaptive
half has a working remeshing loop but no production run.

One claim in this section was \emph{withdrawn} by its own measurement: the
velocity drop is not systematically larger in VOF. See
Section~\ref{sec:dropdef}.}

\subsection{Case definition}

The Newtonian case of Section~\ref{sec:newtonian} with an Oldroyd-B medium
surrounding a Newtonian bubble, following the reference
thesis~\cite{silva-thesis}. Two changes to the geometry and the schedule, both
with a reason:

\begin{itemize}
\item the tank is raised to $[0,1]\times[0,1]\times[0,3]$, so the top and bottom
walls do not influence the polymeric tensor;
\item the final time goes from $t=3$ to $t=4$, because the elastic stress takes
longer than the velocity to reach a steady state.
\end{itemize}

\begin{table}[htbp]
\centering
\caption{Viscoelastic matrix: relaxation time and solvent fraction.}
\label{tab:vematrix}
\small
\begin{adjustbox}{max width=\linewidth}
\pgfplotstabletypeset[
  col sep=space, string type,
  columns/lambda/.style={column name={$\lambda_1$}, column type={c}},
  columns/beta/.style={column name={$\beta$}, column type={c}},
  columns/polimero/.style={column name={Polymer}, column type={c}},
  columns/esperado/.style={column name={Expected shape at $t=4$}, column type={l}},
  every head row/.style={before row=\toprule, after row=\midrule},
  every last row/.style={after row=\bottomrule},
]{\dat{ve-matriz.dat}}
\end{adjustbox}
\end{table}

\subsection{The cusp as the discriminating case}

At $\lambda_1 = 0.4$ and $\beta = 0.2$ the reference reports a \textbf{cusp}
forming at the trailing edge of the bubble: the rising motion deforms the
viscoelastic medium, which exerts elastic force back on the interface, and that
force stretches the rear of the bubble into a point.

This is the case that gives the comparison content rather than making it a tie.
A cusp is a region of high curvature concentrated on a scale that approaches the
cell size, and the two representations should handle it differently: an explicit
surface with surgery can refine locally along the interface, whereas a volume
fraction with a per-cell planar reconstruction cannot represent curvature below
one cell.

\status{The VOF half is complete at the reference resolution --- 2001 steps to
$t=4$, 161 extracted surfaces, all closed, $0.45\%$ volume drift. The
front-tracking half is running with the filter of Section~\ref{sec:tsur}, which
is what makes it comparable at all (Table~\ref{tab:undtime}).}

What the VOF half settles on its own is narrower than the comparison, but it is
not nothing: at $h=1/50$ the rear of the bubble reaches $z_{\min} = 1.805$ with
the centroid at $2.042$, against $1.835$ and $2.054$ at $h=1/25$ --- the trailing
edge is drawn out by $0.018$ more in units of $L$ relative to the centroid when
the mesh is halved. That is the direction a cusp would produce, and it is also
the direction a more resolved planar reconstruction would produce without any
cusp. \missing{the front-tracking half, which is the only one of the two that
can represent curvature below the cell and so the only one that can distinguish
these}

\subsection{A model limit, not a numerical one}

Oldroyd-B has no finite extensibility limit (Section~\ref{sec:formulation}), so
under sustained extension the stress grows without bound. In a preliminary run
with deliberately wrong parameters --- gravity a hundredfold too strong and
surface tension 24 times too weak --- the trace grew monotonically from 0 to 416
in 160 steps with no sign of saturation. With the correct parameters the same
configuration gives \num{3.6e-2} after seven steps, growing slowly. The
distinction matters for reading the viscoelastic results: unbounded growth is
what the model does in extension, and the question is whether it saturates on
the time scale of the rise, not whether it appears.

\subsection{Uniform mesh}

All twelve runs completed 2001 steps to $t=4$ with zero non-finite values.
That is worth stating before any number: Oldroyd-B has no finite extensibility
limit, and the question going in was whether the stress would saturate on the
time scale of the rise or run away.

\begin{table}[htbp]
\centering
\caption{The measured matrix at $t=4$, $25\times25\times75$. ``Drop'' is the
fall in rise velocity after its maximum, relative to that maximum.}
\label{tab:vemeasured}
\small
\begin{adjustbox}{max width=\linewidth}
\pgfplotstabletypeset[
  col sep=space, string type,
  columns={caso,lambda,beta,zc,wmax,tmax,west,queda},
  columns/caso/.style={column name={Run}, column type={l}},
  columns/lambda/.style={column name={$\lambda_1$}, column type={c}},
  columns/beta/.style={column name={$\beta$}, column type={c}},
  columns/zc/.style={column name={$z_c$}, column type={r}},
  columns/wmax/.style={column name={$w_{\max}$}, column type={r}},
  columns/tmax/.style={column name={at $t$}, column type={r}},
  columns/west/.style={column name={$w_{\mathrm{steady}}$}, column type={r}},
  columns/queda/.style={column name={Drop (\%)}, column type={r}},
  every head row/.style={before row=\toprule, after row=\midrule},
  every last row/.style={after row=\bottomrule},
]{\dat{ve-medido.dat}}
\end{adjustbox}
\end{table}

\subsubsection{Against the published numbers}

The reference thesis publishes curves rather than tables, but its text quotes
values, and those are what make a numerical comparison possible at all. They
come with ``approximately'' and two significant figures, which bounds how tight
the comparison can be.

\begin{table}[htbp]
\centering
\caption{Measured against the values the reference thesis states in its text.}
\label{tab:vethesis}
\small
\begin{adjustbox}{max width=\linewidth}
\pgfplotstabletypeset[
  col sep=space, string type,
  columns/caso/.style={column name={Run}, column type={l}},
  columns/zc/.style={column name={$z_c$}, column type={r}},
  columns/zcrel/.style={column name={}, column type={r}},
  columns/wmax/.style={column name={$w_{\max}$}, column type={r}},
  columns/wmaxrel/.style={column name={}, column type={r}},
  columns/west/.style={column name={$w_{\mathrm{steady}}$}, column type={r}},
  columns/westrel/.style={column name={}, column type={r}},
  every head row/.style={before row=\toprule, after row=\midrule},
  every last row/.style={after row=\bottomrule},
]{\dat{ve-contra-tese.dat}}
\end{adjustbox}
\end{table}

Every entry agrees to within 7\%, and most to within 3\%. The signs are not
random, and the next subsection returns to that once the Newtonian reference
puts nine comparisons on the table rather than six.

\subsubsection{The trends, all twelve confirmed}

The thesis makes three ordering claims: more polymer rises higher, reaches a
higher peak velocity, and settles at a higher steady velocity. Across both
methods and both relaxation times, \textbf{all twelve orderings hold}.

\begin{center}
\small
\begin{tabular}{llccc}
\toprule
Method & $\lambda_1$ & $z_c$ ($\beta = 0.8, 0.5, 0.2$) & $w_{\max}$ & $w_{\mathrm{steady}}$ \\
\midrule
FT  & 0.1 & 1.8221, 1.8482, 1.8761 & 0.3642, 0.3698, 0.3754 & 0.3486, 0.3561, 0.3641 \\
FT  & 0.4 & 1.8477, 1.9156, 1.9860 & 0.3714, 0.3896, 0.4092 & 0.3557, 0.3751, 0.3972 \\
VOF & 0.1 & 1.8679, 1.8965, 1.9251 & 0.3942, 0.4002, 0.4041 & 0.3641, 0.3713, 0.3816 \\
VOF & 0.4 & 1.8942, 1.9733, 2.0540 & 0.3933, 0.4158, 0.4280 & 0.3685, 0.3948, 0.4180 \\
\bottomrule
\end{tabular}
\end{center}

\paragraph{One ordering appeared to fail, and the reason is methodological.}
Read before all twelve runs had reached $t=4$, the VOF series at
$\lambda_1 = 0.1$ was \emph{not} monotone in $z_c$: 1.8679, 1.8965, 1.8774. The
third run was at $t = 3.88$ while the second was at $t = 4.00$, and a centre of
mass is an integrated quantity --- comparing it across different end times
compares different things. At $t = 4.00$ the ordering holds. Every table in this
paper that reports an integrated quantity therefore also reports the end time it
was read at.

\subsubsection{The signature the thesis describes, in both methods}

The most specific claim in the reference is not an ordering but a shape: at
$\lambda_1 = 0.4$ and $\beta = 0.2$ the rise velocity falls sharply after its
maximum and then recovers. It appears in both methods --- a drop of $8.1\%$ in
front-tracking and $5.9\%$ in VOF for the cusp case, with recovery in both.

The drop is present in \emph{every} viscoelastic run here. An earlier version of
this section reported it as \emph{systematically larger in VOF}, and that claim
does not survive its own measurement --- the correction is worth stating in full,
because the failure mode is general.

\subsubsection{The asymmetry was in the measurement, not the flow}
\label{sec:dropdef}

The drop had been taken as the fall from the global maximum of $w(t)$ to the
following minimum. On a curve that is nearly flat after the initial overshoot,
that definition is governed by whichever single frame happens to sit highest ---
and in three of the seven VOF runs the ``maximum'' it selects lies at
$t \approx 3.0$--$3.9$, long after the physical overshoot at $t \approx 0.8$.

The two methods do not carry the same frame-to-frame scatter. Locating the
overshoot in the window $t \in [0.3,1.5]$ on a centred five-point mean, and
measuring the scatter of $w$ about that mean:

\begin{table}[htbp]
\centering
\caption{The same drop under two definitions, and the scatter that separates
them. ``Raw'' is global maximum then following minimum; ``robust'' locates the
overshoot physically and smooths.}
\label{tab:dropdef}
\small
\begin{adjustbox}{max width=\linewidth}
\pgfplotstabletypeset[
  col sep=space, string type,
  columns={caso,metodo,lambda,beta,qbruta,tbruta,qrobusta,trobusta,ruido},
  columns/caso/.style={column name={Run}, column type={l}},
  columns/metodo/.style={column name={}, column type={l}},
  columns/lambda/.style={column name={$\lambda_1$}, column type={r}},
  columns/beta/.style={column name={$\beta$}, column type={r}},
  columns/qbruta/.style={column name={raw (\%)}, column type={r}},
  columns/tbruta/.style={column name={$t$ of max}, column type={r}},
  columns/qrobusta/.style={column name={robust (\%)}, column type={r}},
  columns/trobusta/.style={column name={$t$ of max}, column type={r}},
  columns/ruido/.style={column name={scatter (\% of peak)}, column type={r}},
  every head row/.style={before row=\toprule, after row=\midrule},
  every last row/.style={after row=\bottomrule},
]{\dat{queda-robusta.dat}}
\end{adjustbox}
\end{table}

Under the raw definition front-tracking gives 11.5--13.8\% and VOF 11.1--23.1\%,
and VOF looks systematically worse. Under the robust definition front-tracking
gives 5.5--8.1\% (mean 6.8\%) and VOF 5.1--7.0\% (mean 6.1\%): the two
\textbf{agree}, and if anything front-tracking drops slightly more.

What differs between the methods is not the drop but the scatter --- $3.0\%$ of
the peak in front-tracking against $3.8$--$4.9\%$ in VOF, consistent with jitter
in the marching-cubes extraction from which $z_c$ is read. The raw definition
converts that scatter into an apparent physical asymmetry, because a noisier
series offers a higher stray maximum.

So the honest statement is narrower than the one it replaces: \textbf{the
velocity-drop signature the reference describes is present in both methods at
$6\pm1\%$, and these runs do not resolve a difference between them.} The
asymmetry that remains is the one in Section~\ref{sec:newtonian} --- VOF above
the published values on all nine comparisons --- which does not depend on
locating an extremum.

\subsubsection{And it does not change with $h$}

The pending test was whether the drop moves with resolution: a representation
artefact should, the physics should not. The cusp case now exists at both
resolutions, the same quantity computed the same way.

\begin{table}[htbp]
\centering
\caption{The drop at two resolutions, VOF, $\lambda_1=0.4$, $\beta=0.2$.}
\label{tab:dropres}
\small
\begin{adjustbox}{max width=\linewidth}
\pgfplotstabletypeset[
  col sep=space, string type,
  columns/malha/.style={column name={Mesh}, column type={l}},
  columns/h/.style={column name={$h$}, column type={r}},
  columns/queda/.style={column name={drop (\%)}, column type={r}},
  columns/pico/.style={column name={$w_{\max}$}, column type={r}},
  columns/tpico/.style={column name={$t_{\max}$}, column type={r}},
  columns/zcfinal/.style={column name={$z_c(4)$}, column type={r}},
  columns/psifinal/.style={column name={$\Psi(4)$}, column type={r}},
  columns/ruido/.style={column name={scatter (\%)}, column type={r}},
  every head row/.style={before row=\toprule, after row=\midrule},
  every last row/.style={after row=\bottomrule},
]{\dat{queda-resolucao.dat}}
\end{adjustbox}
\end{table}

Halving $h$ moves the drop from $5.9\%$ to $6.6\%$ --- $0.7$ percentage points,
against a frame scatter of $3$--$4\%$ of the peak. The measurement therefore
\emph{bounds} any resolution dependence below its own noise rather than
resolving one, which is consistent with the drop being physical and is not proof
of it. The rise itself converges in the usual way: $z_c(4)$ moves $0.60\%$ and
the final sphericity $0.0046$.

\subsubsection{The shapes}

Figure~\ref{fig:veshape} is the counterpart of the reference's shape figures:
the vertical section through the bubble centre, at $t=4$, for all six
$(\lambda_1,\beta)$ combinations, with the two methods overlaid. As in
Section~\ref{sec:newtonian}, contours are centred on their own centroid so that
what is compared is shape and not rise.

\begin{figure}[htbp]
\centering
\begin{tikzpicture}
\begin{axis}[width=0.325\linewidth, height=0.30\linewidth,
             axis equal image, title={$\lambda_1{=}0.1$, $\beta{=}0.8$}, grid=both, no markers,
             legend pos=south east, legend cell align=left,
             font=\scriptsize, xticklabel pos=bottom]
  \addplot+[thick] table[x=x, y=z] {\dat{secoes/ve-vof-l01-b08-c.dat}};
    \addlegendentry{VOF}
  \addplot+[thick, dashed] table[x=x, y=z] {\dat{secoes/ve-ft-l01-b08-c.dat}};
    \addlegendentry{FT}
\end{axis}
\end{tikzpicture}
\begin{tikzpicture}
\begin{axis}[width=0.325\linewidth, height=0.30\linewidth,
             axis equal image, title={$\lambda_1{=}0.1$, $\beta{=}0.5$}, grid=both, no markers,
             font=\scriptsize, xticklabel pos=bottom]
  \addplot+[thick] table[x=x, y=z] {\dat{secoes/ve-vof-l01-b05-c.dat}};
  \addplot+[thick, dashed] table[x=x, y=z] {\dat{secoes/ve-ft-l01-b05-c.dat}};
\end{axis}
\end{tikzpicture}
\begin{tikzpicture}
\begin{axis}[width=0.325\linewidth, height=0.30\linewidth,
             axis equal image, title={$\lambda_1{=}0.1$, $\beta{=}0.2$}, grid=both, no markers,
             font=\scriptsize, xticklabel pos=bottom]
  \addplot+[thick] table[x=x, y=z] {\dat{secoes/ve-vof-l01-b02-c.dat}};
  \addplot+[thick, dashed] table[x=x, y=z] {\dat{secoes/ve-ft-l01-b02-c.dat}};
\end{axis}
\end{tikzpicture}

\begin{tikzpicture}
\begin{axis}[width=0.325\linewidth, height=0.30\linewidth,
             axis equal image, title={$\lambda_1{=}0.4$, $\beta{=}0.8$}, grid=both, no markers,
             font=\scriptsize, xticklabel pos=bottom]
  \addplot+[thick] table[x=x, y=z] {\dat{secoes/ve-vof-l04-b08-c.dat}};
  \addplot+[thick, dashed] table[x=x, y=z] {\dat{secoes/ve-ft-l04-b08-c.dat}};
\end{axis}
\end{tikzpicture}
\begin{tikzpicture}
\begin{axis}[width=0.325\linewidth, height=0.30\linewidth,
             axis equal image, title={$\lambda_1{=}0.4$, $\beta{=}0.5$}, grid=both, no markers,
             font=\scriptsize, xticklabel pos=bottom]
  \addplot+[thick] table[x=x, y=z] {\dat{secoes/ve-vof-l04-b05-c.dat}};
  \addplot+[thick, dashed] table[x=x, y=z] {\dat{secoes/ve-ft-l04-b05-c.dat}};
\end{axis}
\end{tikzpicture}
\begin{tikzpicture}
\begin{axis}[width=0.325\linewidth, height=0.30\linewidth,
             axis equal image, title={$\lambda_1{=}0.4$, $\beta{=}0.2$}, grid=both, no markers,
             font=\scriptsize, xticklabel pos=bottom]
  \addplot+[thick] table[x=x, y=z] {\dat{secoes/ve-vof-l04-b02-c.dat}};
  \addplot+[thick, dashed] table[x=x, y=z] {\dat{secoes/ve-ft-l04-b02-c.dat}};
\end{axis}
\end{tikzpicture}
\caption{Viscoelastic bubble at $t=4$: vertical section through the centre, VOF
(solid) against front-tracking (dashed). Top row $\lambda_1 = 0.1$, bottom row
$\lambda_1 = 0.4$; polymer content increases to the right ($\beta$ falls).
Contours centred on their own centroid.}
\label{fig:veshape}
\end{figure}

Two things the figure shows that the tables cannot. The first is that the two
methods agree on shape across the whole matrix, including at the highest polymer
content, where the elastic stress is largest. The second is the trailing edge:
at $\lambda_1 = 0.4$, $\beta = 0.2$ --- the bottom-right panel --- the rear of
the bubble is drawn out relative to every other case, which is the cusp the
reference describes.

\paragraph{What this figure does not settle.} Whether the drawn-out rear is
\emph{resolved} or merely \emph{suggested} at this mesh. A cusp is curvature
concentrated on a scale that approaches the cell size, and this is the coarse
mesh. \missing{the reference resolution for this case is running; the comparison
of the two will say whether the feature sharpens with $h$ --- which it should, if
it is the cusp --- or stays at the same width, which would mean the mesh is
setting its scale}

\subsubsection{The Newtonian reference, in the same geometry}

One claim of the reference could not be tested until now: that the viscoelastic
bubble rises higher than a Newtonian one. The Newtonian runs of
Section~\ref{sec:newtonian} cannot serve as that reference --- they use a
different box and end time, and the centre of mass is an \emph{integrated}
quantity, so it is not comparable across either. A Newtonian pair was therefore
run in the thesis geometry: box $1\times1\times3$, to $t=4$.

\begin{table}[htbp]
\centering
\caption{The trend claim, with a reference in the same geometry. Centre of mass
at $t=4$.}
\label{tab:vetendency}
\small
\pgfplotstabletypeset[
  col sep=space, string type,
  columns/caso/.style={column name={Case ($\beta=0.2$)}, column type={l}},
  columns/ve/.style={column name={viscoelastic}, column type={r}},
  columns/newt/.style={column name={Newtonian}, column type={r}},
  columns/veredito/.style={column name={}, column type={l}},
  every head row/.style={before row=\toprule, after row=\midrule},
  every last row/.style={after row=\bottomrule},
]{\dat{ve-tendencia.dat}}
\end{table}

All four hold. And the Newtonian reference itself lands close to the published
value: front-tracking gives $z_c = 1.8055$ against $\approx 1.8$, and a peak
velocity of $0.3603$ against $\approx 0.36$ --- $0.3\%$ and $0.1\%$.

\paragraph{The asymmetry in the signs is now worth stating as a result.}
With the Newtonian reference included there are nine comparisons against
published values per method. \textbf{VOF is above the reference on all nine};
front-tracking is at or below on seven, and its two exceptions are $+0.1\%$ and
$+0.3\%$. The mean absolute deviation is $1.2\%$ for front-tracking against
$3.6\%$ for VOF.

That is no longer the scatter of two or three readings, and it is the opposite
of what the provenance would suggest: the reference was computed with VOF. Two
readings are available and these runs do not separate them --- our
volume-of-fluid differs from theirs in some detail of the reconstruction or the
initial fraction, or the marching-cubes extraction from which our shape
quantities are read is itself biased upward. \missing{the second is testable:
the field-based centroid probe agrees with the surface-based one to
\num{1e-4} in the Newtonian case of Section~\ref{sec:newtonian}, which argues
against extraction bias, but that check has not been repeated here}

\subsection{Adaptive mesh}
\missing{blocked on Section~\ref{sec:amr}}
